\documentclass[10pt]{article}
\usepackage[T1]{fontenc}
\usepackage[utf8]{inputenc}
\usepackage[french]{babel}
\usepackage{amsmath,amssymb,mathrsfs,mathpazo}
\usepackage{geometry}
\geometry{a4paper,margin=2cm}
\usepackage{graphicx,array,tabularx,tabularray,colortbl,multirow,multicol,booktabs,enumitem,pifont,ulem,fancybox,diagbox,cancel}
\usepackage{pstricks,pstricks-add,pst-plot,pst-tree,pst-node,pst-text,pst-func,pst-math,pst-bspline,pst-3dplot,pst-eucl,pst-solides3d}
\let\pscalculate\undefined
\usepackage{pst-bezier}
\usepackage{tikz,pgfplots,tkz-tab}
\usetikzlibrary{arrows,arrows.meta,patterns,calc,positioning,trees,decorations.markings,shapes.geometric,shapes.misc,shapes.arrows,intersections,decorations.pathmorphing,decorations.pathreplacing,backgrounds,shadows,babel}
\pgfplotsset{compat=1.18}
\usepgfplotslibrary{fillbetween}
\usepackage[np]{numprint}
\usepackage{listings,makecell,pas-tableur,twoopt}
\let\barre\undefined
\usepackage[pstricks]{tablvar}
\definecolor{darkgreen}{RGB}{0,100,0}
\providecommand{\red}{\color{red}}
\providecommand{\blue}{\color{blue}}
\providecommand{\green}{\color{green}}
\providecommand{\black}{\color{black}}
\providecommand{\white}{\color{white}}
\providecommand{\gray}{\color{gray}}
\providecommand{\cyan}{\color{cyan}}
\providecommand{\magenta}{\color{magenta}}
\providecommand{\yellow}{\color{yellow}}
\providecommand{\violet}{\color{violet}}
\providecommand{\orange}{\color{orange}}
\providecommand{\euro}{\texteuro}
\providecommand{\EUR}{\texteuro}
\providecommand{\ieme}{\textsuperscript{e}}
\providecommand{\degres}{\ensuremath{^\circ}}
\providecommand{\degree}{\ensuremath{^\circ}}
\providecommand{\vv}[1]{\overrightarrow{#1}}
\providecommand{\Coord}[1]{#1}
\providecommand{\decosix}{}
\providecommand{\starredbullet}{\ensuremath{\bullet}}
\pagestyle{empty}
\setlength{\parindent}{0pt}
\usepackage{adjustbox}\DeclareUnicodeCharacter{2236}{:}\DeclareUnicodeCharacter{2190}{\ensuremath{\leftarrow}}
\providecommand{\R}{\mathbb{R}}
\providecommand{\N}{\mathbb{N}}
\providecommand{\Z}{\mathbb{Z}}
\providecommand{\Q}{\mathbb{Q}}
\providecommand{\D}{\mathbb{D}}
\providecommand{\C}{\mathbb{C}}
\providecommand{\e}{{\mathrm{e}}}
\providecommand{\vect}[1]{\overrightarrow{#1}}
\providecommand{\vectt}[1]{\overrightarrow{\mathrm{#1}}}
\providecommand{\vvt}[1]{\overrightarrow{\mathrm{#1}}}
\providecommand{\np}[2][]{{#2~#1}}
\providecommand{\vv}[1]{{\overrightarrow{#1}}}
\renewcommand{\d}{{\mathrm{d}}}
\providecommand{\degres}{{^\circ}}
\providecommand{\degre}{{^\circ}}
\providecommand{\blue}{\color{blue}}
\providecommand{\red}{\color{red}}
\providecommand{\enskip}{\quad}
\providecommand{\llbracket}{\lbrack\!\lbrack}
\providecommand{\rrbracket}{\rbrack\!\rbrack}
\providecommand{\euro}{{\text{€}}}
\providecommand{\no}{{\text{n°}}}
\providecommand{\interval}[2]{{[#1;#2]}}
\providecommand{\intervFO}[2]{{[#1;#2[}}
\providecommand{\dotfill}{\ldots}

\providecommand{\ds}{\displaystyle}
\providecommand{\textding}[1]{#1}
\providecommand{\Oij}{\(\mathrm{O};\vec{\imath},\vec{\jmath}\)}
\providecommand{\Oijk}{\(\mathrm{O};\vec{\imath},\vec{\jmath},\vec{k}\)}
\providecommand{\Ouv}{\(\mathrm{O};\vec{u},\vec{v}\)}
\providecommand{\intervFF}[2]{\ensuremath{[#1;#2]}}
\renewcommand{\euro}{\ifmmode\text{\texteuro}\else\texteuro\fi}
\DeclareUnicodeCharacter{00B0}{\ifmmode^{\circ}\else\textdegree\fi}
\DeclareUnicodeCharacter{20AC}{\ifmmode\text{\texteuro}\else\texteuro\fi}
% Énoncés : Éducation nationale. Transcription et corrigés : APMEP, Denis Vergès et contributeurs. Découpage et classement : CoopMaths, Sébastien Lozano, Éric Elter et contributeurs. Adaptation de la présentation : Maths Explorer.
% CC BY-SA : https://coopmaths.fr/www/about/
% https://coopmaths.fr/alea/static/bac/2023/tex/bac_2023_03_sujet1_etrangers_groupe2_1_cor.tex
\begin{document}
\noindent\textit{APMEP / CoopMaths. CC BY-SA. Présentation adaptée par Maths Explorer.}\par\medskip
\medskip

\textbf{Partie A}

\medskip

On considère la fonction $g$ définie sur l'intervalle $]0~;~+\infty[$ par
\[g(x) = \ln \left(x^2\right) + x - 2\]
\begin{enumerate}
\item ~
$\bullet~~$On a $\displaystyle\lim_{x \to 0} \ln x^2 = - \infty$ et $\displaystyle\lim_{x \to 0} x = 0$, d'où par somme de limites : $\displaystyle\lim_{x \to 0} g(x) = - \infty$.

$\bullet~~$$\displaystyle\lim_{x \to + \infty} \ln x^2  = + \infty$ et $\displaystyle\lim_{x \to + \infty} x  = + \infty$, d'où  par somme de limites : $\displaystyle\lim_{x \to + \infty} g(x) = + \infty$.
\item 


Sur $]0~;~+\infty[$, on a $g(x) = 2\ln x + x - 2$, d'où :

$g'(x) = 2\times \dfrac1x + 1$.
On a par composition de la dérivation :

$g'(x) = 2 \times \dfrac1x + 1 = \dfrac{2 + x}{x}$.

Comme $x > 0$, le signe de $g'(x)$ est celui de $2 + x$.

Or $x > 0 \implies 2 + x > 2 > 0$, donc $g'(x) > 0$ sur l'intervalle $]0~;~+\infty[$.

La fonction $g$ est donc croissante et d'après la question précédente de moins l'infini à plus l'infini.
\item \begin{enumerate}
\item 
D'après la question précédente $g$ est continue car dérivable sur $]0~;~+\infty[$ et croit de moins l'infini à plus l'infini : d'après le théorème des valeurs intermédiaires il existe donc un réel unique $\alpha \in ]0~;~+\infty[$ tel que $g(\alpha) = 0$.
\item ~
La calculatrice donne :

$g(1) = - 1$ et $g(2) \approx1,4$, donc $1 < \alpha < 2$ ;

$g(1,3) \approx - 0,175$ et $g(1,4) \approx 0,07$, donc $1,3 < \alpha < 1,4$ ;

$g(1,37) \approx - 0,0004$ et $g(1,38) \approx 0,02$, donc $1,37 < \alpha < 1,38$
\end{enumerate}
\item D'où le tableau de signe :


\begin{center}
\psset{unit=1cm}
\begin{pspicture}(7,1)
\psframe(7,1)\psline(0,0.5)(7,0.5)\psline(1,0)(1,1)
\uput[u](0.5,0.4){$x$}\uput[u](1.15,0.4){0}\uput[u](4,0.4){$\alpha$}\uput[u](6.5,0.4){$+ \infty$}
\uput[u](0.5,-0.1){$g(x)$}\uput[u](2.5,-0.1){$-$}\uput[u](4,-0.1){$0$}\uput[u](5.5,-0.1){$+$}
\end{pspicture}
\end{center}
\end{enumerate}

\medskip

\textbf{Partie B}

\medskip

On considère la fonction $f$ définie sur l'intervalle $]0~;~+\infty[$ par:
\[f(x) = \dfrac{(x-2)}{x}\ln(x)\]

\begin{enumerate}
\item 
	\begin{enumerate}
		\item 
On a successivement :

$\displaystyle\lim_{x \to 0} (x - 2) = - 2$ et $\displaystyle\lim_{x \to 0} x  = 0$, d'où par quotient de limites $\displaystyle\lim_{x \to 0} \dfrac{(x - 2)}{x} = - \infty$ ;

D'autre part $\displaystyle\lim_{x \to 0} \ln x = - \infty$, donc par produit de limites : $\displaystyle\lim_{x \to 0} f(x) = + \infty.$
		\item 
		Géométriquement le résultat précédent montre que l'axe des ordonnées est asymptote verticale à la courbe $\mathcal{C}_f$.
	\end{enumerate}
\item 
$\bullet~~$On a $\dfrac{x - 2}{x} = \dfrac{x\left(1 - \frac2x\right)}{x} = 1 - \frac2x$

De $\displaystyle\lim_{x \to + \infty} \frac2x = 0$, il suit $\displaystyle\lim_{x \to + \infty}\dfrac{x - 2}{x} = 1$.

$\bullet~~$D'autre part on sait que $\displaystyle\lim_{x \to + \infty}\ln x = + \infty$. 

Donc par produit de limites, on a $\displaystyle\lim_{x \to + \infty}f(x) = + \infty$.
\item On admet que la fonction $f$ est dérivable sur l'intervalle $]0~;~+\infty[$,


Sur l'intervalle $]0~;~+\infty[$, on a par dérivation du produit :

$f'(x) = \dfrac{1 \times x - 1 \times (x - 2)}{x^2}\ln x + \dfrac{x - 2}{x} \times \dfrac1x = \dfrac{x - x + 2}{x^2}\ln x + \dfrac{x - 2}{x^2} = \dfrac{2\ln x + x - 2}{x^2} = \dfrac{g(x)}{x^2}$.
\item  
Comme $x^2 > 0$ sur $]0~;~+\infty[$, le signe de $f'(x)$ sur cet intervalle est celui de $g(x)$.

On a vu à la question A. 4. que :

$\bullet~~$ :sur l'intervalle $]0~;~\alpha[, \: g(x) < 0$ : la fonction $f$ est donc décroissante sur $]0~;~\alpha[$ ;

$\bullet~~$ :sur l'intervalle $]\alpha~;~+ \infty[, \: g(x) > 0$ : la fonction $f$ est donc croissante sur $]\alpha~;~\infty[$.
\end{enumerate}
\medskip
\textbf{Partie C}
\medskip


\[d(x) = f(x) - \ln x = \dfrac{(x-2)}{x}\ln(x) - \ln x = \ln x\left[\dfrac{(x-2)}{x} - 1\right] = \ln x\left[\dfrac{x - 2 - x}{x}\right] = \dfrac{- 2}{x}\ln x.\]

La position relative de la courbe $\mathcal{C}_f$ et de la courbe représentative de la fonction $\ln$ est donnée par le signe de la fonction $d$.

Comme $x > 0$, le signe de $d(x)$ est celui du produit $- 2\ln x$.

On dresse donc un tableau de signes :

\begin{center}
\psset{unit=1cm}
\begin{pspicture}(8,2)
\psframe(8,2)\psline(0,0.5)(8,0.5)\psline(2,0)(2,2)\psline(0,1)(8,1)\psline(0,1.5)(8,1.5)
\psline(2.1,0)(2.1,1.5)\psline(2.2,0)(2.2,1.5)
\uput[u](1,1.4){$x$}\uput[u](2.15,1.4){0}\uput[u](5,1.4){1}\uput[u](7.5,1.4){$+ \infty$}
\uput[u](1,0.9){$- 2$}\uput[u](3.5,0.9){$-$}\uput[u](6.5,0.9){$-$}
\uput[u](1,0.4){$\ln x$}\uput[u](3.5,0.4){$-$}\uput[u](5,0.4){$0$}\uput[u](6.5,0.4){$+$}
\uput[u](1,-0.1){$- 2\ln x$}\uput[u](3.5,-0.1){$+$}\uput[u](5,-0.1){$0$}\uput[u](6.5,-0.1){$-$}

\end{pspicture}
\end{center}
Conclusion : sur l'intervalle $]0~;~1[$, la courbe $\mathcal{C}_f$ est au dessus de la courbe logarithme népérien et,

sur l'intervalle $]1~;~+ \infty[$ la courbe $\mathcal{C}_f$ est au dessous de la courbe logarithme népérien.

\begin{center}
\psset{unit=1cm,arrowsize=2pt 3}
\begin{pspicture*}(-0.5,-2)(7,2)
\psaxes[linewidth=1.25pt]{->}(0,0)(0,-2)(7,2)
\psgrid(0,-2)(7,5)
\psplot[plotpoints=2000,linecolor=red,linewidth=1.25pt]{0.1}{7}{x 2 sub x div x ln mul}\uput[dr](6,1.2){\red$\mathcal{C}$}
\psplot[plotpoints=2000,linecolor=blue,linewidth=1.25pt]{0.1}{7}{x ln}\uput[u](4,1.4){\blue $y = \ln x$}
\end{pspicture*}
\end{center}
\bigskip

\end{document}
